\section{Matter and the Forces}
\label{sec:matter}

Matter and the forces are patterns on the substrate. This section walks the field content and the
gauge sector (P7, P8, P18 to P25, P55).

\paragraph{Particles as persistent patterns (P7).}
A particle is not a point but a stable knot, a pattern that holds its shape as the mesh evolves
around it, like a vortex in a fluid. Different particles are different stable knots. Spin is a
feature of the knot that cannot be smoothly undone, and so is quantised. Mass is a gap, the least
energy to excite the pattern, and a massive pattern obeys the relativistic energy-momentum relation.

\paragraph{The field content by spin (P18 to P22).}
The four field types sorted by spin all appear as patterns: the scalar, the matter spinor, the
light-carrying vector, and the tensor of gravity. The dark sector and the full field content are
covered in P18 to P22. The ternary tone realises the spinor structure (the two-valued plus neutral
character of a fermionic degree of freedom).

\paragraph{Gauge fields and confinement (P8).}
A force is the twist a tone picks up crossing a note, the Wilson construction \citep{wilson1974confinement} that
lattice gauge theory already uses. P8 reproduces confinement, the reason quarks are never seen
alone, together with the index theorem relating the topology of the gauge field to the fermion zero
modes \citep{neuberger1998exactly}, and a chiral condensate. The chiral gauge theory itself is treated in
Section~\ref{sec:standardmodel}.

\paragraph{Operators from the action (P23, P55).}
The operators for light and for gravity are not put in by hand. P23 derives the gauge and graviton
operators as the gentle-vibration (small-fluctuation) limit of one action on the mesh, and P55
shows all the bosonic sectors coming from a single operator on a single mesh, a genuine unification
at the level of the substrate.

\paragraph{Electroweak breaking and mass (P25).}
P25 covers electroweak symmetry breaking: a background scalar settles to a nonzero value, a
Higgs-like mechanism, and the gauge bosons and fermions acquire mass from their coupling to it. Mass
is generated, not inserted. The pattern of fermion masses, the hierarchy, is taken up in
Section~\ref{sec:standardmodel}.
